% REV02: proofs brought together in 02_dualidad_t.tex and 03_pantallas_coordinadas.tex.
\subsection{Exceptional incidence and coordination of the screens}
\label{exc:pantallas-dualidad}

The direction selected by \(K\) and the Leech lattice arise
from different representations of the constructed incidence. The former
determines the orthogonal reduction \(12\to11\to10\); the latter
enters the isotropic reduction \(26\to24\). Their coordination
preserves the preceding space, its relations and the information conveyed
by each projection.

The sections on T-duality and coordinated screens in \cite{hmtX} develop
this coordination jointly. The former brings together the involution of
integer coordinates, unitary conjugation with its domains and the
change of residual section retaining the quotient memory. The latter
constructs the morphism \(\mathcal R_M^{\mathrm{alg}}\)
defined there and proves the isomorphism of the
quotient presentations in the corresponding theorem on isomorphic screens.
The graph preserves the coordinate in \(V_{11}\) together with its projection
in \(V_{10}\), while the isotropic quotient recovers the
Leech lattice. The respective kernels and codomains are thereby preserved.

The action ellipse is constructed before duality is applied and fixes
its dimensionless radius together with the reciprocal radius
(the action-ellipse section of \cite{hmtX}). The compact sector then connects
exceptional incidence with a determined quadratic form. The
fields, action and operators of the subsequent physical realization
are presented in their own domain, preserving this algebraic
antecedent and its internal proofs.
